% ============================================================
%  MANUSCRIPT PLAN: how_cells_decide_NESM.tex
%  Stochastic Thermodynamics of Single-Cell Commitment:
%  A Non-Equilibrium Analysis of B. subtilis Sporulation
% ============================================================
%
%  Target journals: Physical Review X / eLife (physics) / PRE
%  Framing: NESM + SHPN as the primary language; Fujita biology
%           as validation prerequisite, not main result.
%
%  Session: 2026-07-08
%  Status: PLAN — skeleton only, content to be filled part by part.
% ============================================================

% ------------------------------------------------------------
%  FRAMING SHIFT vs. plos_one_how_cells_decide.tex
% ------------------------------------------------------------
%
%  Current (PLOS ONE):
%    Main question  — Can SHPN reproduce Fujita biology?
%    Thermodynamics — A "framework section" before Results
%    Abstract lead  — "how cells decide"
%    Discussion     — Biological implications
%
%  NESM version:
%    Main question  — What does NESM tell us about single-cell commitment?
%    Biology        — Validation prerequisite (Results I, brief)
%    Abstract lead  — "stochastic thermodynamics of commitment"
%    Discussion     — What gamma != MF means physically; generality
%
% ------------------------------------------------------------
%  PROPOSED STRUCTURE
% ------------------------------------------------------------

% 1. ABSTRACT (~ 250 words)
%    S1: Gap — cells make irreversible commitment decisions;
%              no thermodynamic characterisation exists at single-cell level.
%    S2: Framework — SHPN as computational substrate for stochastic
%              thermodynamics; five quantities derivable from topology alone.
%    S3: Validation — B. subtilis sporulation reproduces Fujita & Losick
%              (2005) Q1/Q2/Q3 constraints.
%    S4: Main finding — gamma_sub = 0.150 +/- 0.044 (non-MF, non-mean-field
%              universality class); N_c = 1346 uM; EPR cost concentrated at
%              the sensing step (KinA 97.5%).
%    S5: Generality — framework applies to any commitment circuit with
%              known topology.

% 2. INTRODUCTION (~ 600 words, 4 paragraphs)
%    P1: The thermodynamic gap.
%        Cells commit irreversibly (differentiation, sporulation, apoptosis).
%        Classical thermodynamics cannot describe small-N, NESS systems.
%        Stochastic thermodynamics (Seifert 2012) provides exact trajectory-
%        level thermodynamics for Markov chains — but has never been
%        systematically applied to regulatory commitment circuits.
%    P2: Why a unified model is needed.
%        Regulatory decisions involve both high-copy continuous (ODE) and
%        low-copy discrete stochastic regimes coexisting in the same network.
%        Separate formalisms (ODE vs. Boolean vs. CTMC) cannot address this.
%        SHPN (Simao 2025) unifies them in a single 13-tuple formalism with
%        a provable firing semantics.
%    P3: B. subtilis sporulation as working model.
%        Binary, irreversible, well-characterised (Fujita & Losick 2005).
%        Physiological stochasticity is the mechanism (bet-hedging).
%        Three published experimental constraints to validate against.
%        41 places, 36 transitions — tractable but non-trivial.
%    P4: Here we show.
%        We construct a SHPN model, validate against Fujita constraints,
%        then apply stochastic thermodynamics to compute the five quantities.
%        We find gamma != 0.5 (non-mean-field), N_c = 1346 uM, and
%        that 97.5% of the EPR cost of commitment is paid at the sensing step.

% 3. THEORY I — THE SINGLE CELL AS AN OPEN NON-EQUILIBRIUM SYSTEM
%    (~ 800 words)
%
%    3.1 Non-equilibrium steady states and the cell
%        - A cell maintains steep gradients, pumps ions, consumes ATP.
%        - It operates at a NESS, not equilibrium.
%        - Classical thermodynamics (Gibbs, Boltzmann) applies only at
%          equilibrium; NESM is the correct framework.
%        - Key references: Schnakenberg 1976, Hill 1977, Seifert 2012,
%          Rao & Esposito 2018.
%
%    3.2 Stochastic thermodynamics: entropy production along a trajectory
%        - System described by a CTMC with master equation.
%        - Three conditions for applicability:
%          (i)  CTMC (Master equation)
%          (ii) Local detailed balance: k+/k- = exp(Delta_mu / k_B T)
%          (iii) Explicit driving reservoirs
%        - Trajectory-level EPR: sigma[omega] = ln P[omega] / P[omega_tilde]
%          (exact, not an approximation, valid for any N)
%        - Second Law: <sigma> >= 0, equality only at equilibrium.
%        - Key result: this is rigorous for ~10 molecules.
%
%    3.3 Subsystem thermodynamics and the Markovian reservoir assumption
%        - Full cell has N >> 1 processes; we model k << N.
%        - Esposito & Van den Broeck (2010): subsystem EPR is separable
%          when coupling to environment is Markovian.
%        - sigma_sub = sigma_total - sigma_rest - sigma_interface
%        - Interface term vanishes in NESS under Markovian coupling.
%        - For the sporulation subnetwork:
%          * ATP pool = chemical bath (k_B T ln([ATP]/[ADP][Pi]) = driving)
%          * INITIAL_NUTRIENTS = nutrient reservoir
%          * Rest of cell = effective Markovian bath
%        - Biological evidence: Fujita validation succeeds without modelling
%          unmodelled processes -> weak coupling confirmed empirically.
%        - Scope boundary: quantities characterise the commitment circuit
%          during the decision window, not full cellular lifecycle.

% 4. THEORY II — SHPN AS COMPUTATIONAL SUBSTRATE FOR SINGLE-CELL NESM
%    (~ 700 words)
%
%    4.1 The 13-tuple Bio-PN formalism
%        SPN = (P, T, F, W, M_0, Phi, C, F_t, Psi, F_s, W_s, lambda, theta)
%        - Firing semantics (brief, forward to Supplementary for full proof)
%        - Hybrid regime: stochastic (tau-leaping) + continuous (ODE)
%          coexist in one model.
%
%    4.2 Why SHPN maps to stochastic thermodynamics
%        - Stochastic transitions: tau-leaping is a Poisson approximation
%          to the Master equation of a CTMC.
%          -> EPR of stochastic part: exact within tau-leaping approximation.
%        - ODE transitions: continuous flow; EPR is deterministic limit
%          of the full stochastic EPR.
%        - The hybrid formalism separates:
%          Layer A (stochastic, low-copy): commitment decision — fluctuations
%                  are the biological mechanism.
%          Layer B (continuous, high-copy): execution programme — ODE regime.
%
%    4.3 Arc semantics as thermodynamic ports
%        - Normal arcs: stoichiometric flows, contribute to EPR.
%        - Signal-flow arcs: regulatory dissipation channels; consumptive,
%          cascade-gated by theta_eff (basin floor = min dissipation for
%          downstream activation).
%        - Test arcs: zero-mass-transfer sensing; EPR = 0 for the arc,
%          but they gate transitions that do dissipate.
%        - Inhibitor arcs: presence-check with threshold theta; EPR = 0
%          for arc but determines transition enabling.
%        -> The arc graph encodes the thermodynamic topology of the circuit.
%
%    4.4 The five topological-thermodynamic quantities
%        Quantity 1: eta_L (topological load factor)
%            eta_L = |A| / (|P| * |T|)  [arc density of commitment layer]
%            Interpretation: informational load of the cascade per
%            place-transition pair. Topology-constant (NET-T1 result).
%        Quantity 2: W/k_B T (dimensionless work ratio)
%            W = integral of J * A dt along trajectory
%            W/kBT >> 1 -> Fluctuation Theorems inapplicable (as found)
%        Quantity 3: N_c (critical nutrient level)
%            Order parameter transition point; found at 1346 uM.
%        Quantity 4: gamma_sub (sub-critical exponent)
%            CV^2_bin ~ |N_0 - N_c|^(-gamma); measures universality class.
%        Quantity 5: EPR decomposition
%            Per-transition dissipation budget; identifies thermodynamic
%            bottleneck of the decision.

% 5. RESULTS I — BIOLOGICAL VALIDATION (prerequisite, ~ 1.5 pages)
%    "The subnetwork model is biologically grounded before thermodynamic
%     analysis. Three independent experimental constraints from Fujita &
%     Losick (2005) are reproduced."
%
%    Fig 2: Q1 — abrupt vs. natural route, sporulation % + ATP trajectory
%    Fig 3: Q2 — sigma_H as irreversible gate (k=0 -> 0%)
%    Fig 4: Q3 — N_c titration, commitment timing, sigma_H peak
%    Table 1: Fujita protocol sweep results

% 6. RESULTS II — THERMODYNAMIC CHARACTERISATION OF COMMITMENT
%    (~ 4 pages)
%
%    6.1 Topological load and commitment modes (NET-T1, Fig 5 left)
%        - eta_L ~ 1.82-1.88e-8, topology-constant over HL range 20-90 min
%        - Two modes: Mode A (passive, short HL) vs Mode B (SinR-gated, WT)
%        - Interpretation: the arc graph constrains the thermodynamic load
%          independently of parameter variation.
%
%    6.2 Free energy landscape and work (NET-T2/T3, Fig 5 right + Fig 6a)
%        - Waddington landscape U_eff(phi,t) = -ln P_ss(phi,t)
%          (rigorously defined, not a metaphor)
%        - W/kBT = 3.72e7: commitment is deep in the non-linear regime
%          -> Fluctuation Theorems (Jarzynski/Crooks) are inapplicable
%          -> This is a result, not a limitation
%        - FWD perturbation (enrich nutrients): aborts sporulation 38%->0%
%        - REV perturbation (starve more): triggers minimal 0%->2%
%          -> Asymmetric response = signature of irreversibility
%
%    6.3 Critical transition and universality class (NET-T4, Fig 6b)
%        - N_c = 1346 uM (crossover from sub- to supra-critical)
%        - Critical slowing down: mean commitment time 317->337 min near N_c
%        - CV^2_bin ~ |N_0 - N_c|^(-gamma_sub), gamma_sub = 0.150 +/- 0.044
%        - Mean-field prediction: gamma_MF = 0.5
%        - Deviation: 19 sigma below MF
%        - Physical interpretation: the phosphorelay cascade feeding into
%          the SinR/SinI gate introduces effective long-range correlations
%          that suppress fluctuation amplification near N_c.
%
%    6.4 Entropy production decomposition (NET-T5, Fig 7)
%        - Total EPR: 5.6 -> 0.9 -> ~0 kBT/s (tracks commitment progress)
%        - KinA dominates stochastic EPR: 97.5%
%          -> The thermodynamic cost of commitment is paid at the SENSING step
%          -> Execution programme (sigma factors, coat synthesis) is cheap
%        - Septation asymmetry: 14.3x higher EPR in sporulating vs veg cells
%          -> Entropy production signature of the symmetry-breaking event
%        - Biological interpretation: bet-hedging is thermodynamically
%          inexpensive in execution but expensive in sensing.

% 7. DISCUSSION (~ 700 words, 4 paragraphs)
%
%    P1: What gamma != MF means physically.
%        - MF assumes uncorrelated fluctuations (each molecule decides
%          independently near the transition).
%        - The phosphorelay (KinA->Spo0F->Spo0A) is a 3-step cascade with
%          positive feedback (Spo0A~P activates its own transcription).
%        - Signal-flow arcs introduce PreemptionCheck coupling between
%          commitment layer transitions -> effective correlation length
%          larger than a single step -> non-MF universality.
%        - This is the first report of a non-MF critical exponent for a
%          biological regulatory circuit computed from its topology.
%
%    P2: The Markovian reservoir assumption — scope and limitations.
%        - Valid when: sporulation timescale >> metabolic correlation time.
%        - Breaks when: ATP depletion is fast relative to KinA response
%          (deep starvation, starved-then-refed experiments).
%        - Verified empirically: Fujita validation succeeds -> weak coupling.
%        - Honest scope: thermodynamics of the commitment decision window,
%          not the full cellular lifecycle.
%
%    P3: Why topology encodes thermodynamics.
%        - eta_L is topology-constant (NET-T1): arc density does not change
%          with half-life parameterisation.
%        - EPR bottleneck at KinA is a structural property: the sensing step
%          (KinA autophosphorylation) is the single irreversible input to the
%          cascade; all downstream steps are thermodynamically downstream of it.
%        - Connection to Thermodynamic Uncertainty Relations (Barato & Seifert
%          2015): precision of the commitment decision is bounded below by the
%          EPR. The bet-hedging window (700-1800 uM) is the regime where
%          precision is deliberately sacrificed for population-level diversity.
%
%    P4: Generality of the framework.
%        - The five quantities are computable for any SHPN model with:
%          (i) known arc topology, (ii) stochastic rates, (iii) explicit
%          chemical reservoirs.
%        - Candidate circuits: Caulobacter cell-cycle (CtrA/DnaA), E. coli
%          SOS response (LexA/RecA), mammalian apoptosis (Bcl-2 family).
%        - The framework does not require new experiments; it extracts
%          thermodynamic information from models already validated against
%          existing data.

% 8. CONCLUSION (~ 150 words)
%    - We have established that single-cell commitment decisions have a
%      computable thermodynamic characterisation within the NESM framework.
%    - SHPN provides the natural computational substrate: its arc semantics
%      map directly onto thermodynamic ports of an open CTMC system.
%    - For B. subtilis sporulation: gamma_sub = 0.150 != MF, N_c = 1346 uM,
%      EPR cost concentrated at the sensing step.
%    - Five topology-derivable quantities constitute a thermodynamic
%      fingerprint of a commitment circuit.
%    - The Markovian reservoir assumption restricts scope to the decision
%      window; full cellular thermodynamics is a future direction.

% 9. SUPPLEMENTARY
%    S1: Full 13-tuple Bio-PN formalism and firing semantics proofs
%    S2: All sweep parameters, model details, v9 arc table
%    S3: Derivation of eta_L topology-constancy condition
%    S4: Proof that W/kBT >> 1 implies FT inapplicability for commitment
%    S5: Critical exponent fitting procedure and uncertainty analysis

% ============================================================
%  KEY REFERENCES
% ============================================================
%  Seifert, U. (2012). Rep. Prog. Phys. 75, 126001.
%  Rao, R. & Esposito, M. (2018). New J. Phys. 20, 023007.
%  Esposito, M. & Van den Broeck, C. (2010). PRL 104, 090601.
%  Barato, A.C. & Seifert, U. (2015). PRL 114, 158101.
%  Schnakenberg, J. (1976). Rev. Mod. Phys. 48, 571.
%  Jarzynski, C. (1997). PRL 78, 2690.
%  Crooks, G. (1999). PRE 60, 2721.
%  Fujita, M. & Losick, R. (2005). Genes Dev. 19, 2236.
%  Simao, E. (2025). [SHPN formalism paper]
% ============================================================
